% =============================================
% File: appendix.tex
% Phụ lục — cấu hình theo chuẩn HCMUTE NCKH
% Dùng \chapter* + \begingroup\fontsize{13pt}{15.6pt}\selectfont
% (giống acknowledgments.tex, tránh \appendix gây lỗi titleformat)
% =============================================


% ============================================================
% Phụ lục A — Mã nguồn bộ điều khiển PI trên ESP32
% ============================================================
\chapter*{PHỤ LỤC A\\[0.4em]MÃ NGUỒN BỘ ĐIỀU KHIỂN PI TRÊN ESP32}
\addcontentsline{toc}{chapter}{Phụ lục A — Mã nguồn bộ điều khiển PI (ESP32)}
\label{app:esp32_code}

\begingroup
\setstretch{1.5}
\fontsize{13pt}{15.6pt}\selectfont
\setlength{\parskip}{6pt}

Đoạn mã dưới đây là cấu trúc chính của chương trình điều khiển PI vận tốc
triển khai trên vi điều khiển ESP32 WROOM-32 (Arduino framework, C++).
Chu kỳ điều khiển $\Delta t = 10\,$ms, giới hạn tích phân $\pm 1000$,
giới hạn PWM $[-255,\,255]$.

\endgroup

\begin{lstlisting}[language=C++, caption={Vòng điều khiển PI vận tốc trên ESP32},
                   label=lst:esp32_pi]
// -- Cau hinh phan cung --
#define ENC_A    18
#define ENC_B    19
#define PWM_PIN  25
#define DIR_PIN1 26
#define DIR_PIN2 27

volatile long encoderCount = 0;

void IRAM_ATTR encoderISR() {
    encoderCount += (digitalRead(ENC_B) == HIGH) ? 1 : -1;
}

// -- Tinh van toc (RPM) --
float VelCalculate(long &posNow, float dt) {
    float vel = (posNow / 200.0f) / dt * 60.0f;
    posNow = 0;
    return vel;
}

// -- Bien PI van toc --
float Kp = 0.1854f;           // tu GA (setpoint hien tai)
float Ki = 5.288f;
float integral  = 0.0f;
const float DT        = 0.01f;
const float INT_LIMIT = 1000.0f;

// -- Vong dieu khien (goi moi 10 ms) --
void piControl(float velReq, float velCur) {
    float err = velReq - velCur;

    integral += err * DT;
    integral  = constrain(integral, -INT_LIMIT, INT_LIMIT); // anti-windup

    float u = Kp * err + Ki * integral;
    u = constrain(u, -255.0f, 255.0f);

    int pwmVal = (int)fabs(u);
    digitalWrite(DIR_PIN1, u >= 0 ? HIGH : LOW);
    digitalWrite(DIR_PIN2, u >= 0 ? LOW  : HIGH);
    ledcWrite(0, pwmVal);
}
\end{lstlisting}

\begingroup
\fontsize{13pt}{15.6pt}\selectfont
\noindent\textit{Ghi chú:} Tham số $K_p$, $K_i$ được cập nhật theo setpoint vận tốc
hiện tại dựa trên bảng tra (Phụ lục~B) hoặc đầu ra suy luận mạng ANN.
\endgroup


% ============================================================
% Phụ lục B — Bảng tham số PI tối ưu tại 9 mức setpoint
% ============================================================
\chapter*{PHỤ LỤC B\\[0.4em]BẢNG THAM SỐ PI TỐI ƯU TẠI 9 MỨC SETPOINT}
\addcontentsline{toc}{chapter}{Phụ lục B — Bảng tham số PI tối ưu (GA)}
\label{app:pi_table}

\begingroup
\setstretch{1.5}
\fontsize{13pt}{15.6pt}\selectfont
\setlength{\parskip}{6pt}

Bảng dưới tổng hợp đầy đủ bộ tham số PI vận tốc tối ưu
$\{K_{pv}^*,\,K_{iv}^*\}$ xác định bằng giải thuật di truyền (GA)
cho 9 mức setpoint từ 250 đến 2250~RPM.
Các tham số vòng ngoài vị trí $(K_{pp}^*, K_{ip}^*, K_{dp}^*)$ giữ
không đổi ở tất cả các mức setpoint.

\begin{table}[H]
\centering
\caption{Bộ tham số PI vận tốc tối ưu (GA) tại 9 mức setpoint}
\label{tab:pi_optimal_full}
\renewcommand{\arraystretch}{1.35}
\begin{tabular}{cccccc}
\toprule
\textbf{Setpoint} & \multirow{2}{*}{$K_{pv}^*$} & \multirow{2}{*}{$K_{iv}^*$}
    & \textbf{Overshoot} & \textbf{$t_s$} & \textbf{SSE} \\
\textbf{(RPM)} & & & \textbf{(\%)} & \textbf{(s)} & \textbf{(RPM)} \\
\midrule
250  & 0{,}1854 & 5{,}288 & 3{,}2 & 0{,}38 & 0{,}8 \\
500  & 0{,}1721 & 5{,}104 & 3{,}8 & 0{,}42 & 1{,}1 \\
750  & 0{,}1698 & 4{,}932 & 4{,}1 & 0{,}45 & 1{,}4 \\
1000 & 0{,}1612 & 4{,}847 & 4{,}3 & 0{,}48 & 1{,}7 \\
1250 & 0{,}1583 & 4{,}721 & 4{,}5 & 0{,}51 & 2{,}0 \\
1500 & 0{,}1547 & 4{,}658 & 4{,}7 & 0{,}53 & 2{,}3 \\
1750 & 0{,}1512 & 4{,}512 & 4{,}8 & 0{,}56 & 2{,}6 \\
2000 & 0{,}1487 & 4{,}389 & 4{,}9 & 0{,}59 & 2{,}9 \\
2250 & 0{,}1463 & 4{,}274 & 5{,}0 & 0{,}62 & 3{,}1 \\
\midrule
\multicolumn{6}{l}{%
  \textbf{Tham số vòng ngoài (cố định):}
  $K_{pp}^* = 0{,}42$;\quad $K_{ip}^* = 0{,}18$;\quad $K_{dp}^* = 0{,}05$%
} \\
\bottomrule
\end{tabular}
\end{table}

\noindent Trong đó: $t_s$ = thời gian xác lập (sai số $< 2\%$);
SSE = sai số xác lập trung bình.
Tất cả chỉ số đều trong ràng buộc thiết kế
(overshoot $\leq 5\%$, $t_s \leq 0{,}65\,$s).

\endgroup


% ============================================================
% Phụ lục C — Động học cơ cấu 5-bar Parallel
% ============================================================
\chapter*{PHỤ LỤC C\\[0.4em]ĐỘNG HỌC CƠ CẤU 5-BAR PARALLEL}
\addcontentsline{toc}{chapter}{Phụ lục C — Động học cơ cấu 5-bar parallel}
\label{app:fivebar}

\begingroup
\setstretch{1.5}
\fontsize{13pt}{15.6pt}\selectfont
\setlength{\parskip}{6pt}

\textbf{Thông số hình học:}
\begin{itemize}[leftmargin=1.5cm]
    \item Khoảng cách hai tâm khớp chủ động: $A_1A_2 = 105\,$mm
    \item Chiều dài khâu chủ động: $L_{act} = 100\,$mm
    \item Chiều dài khâu bị động: $L_{pas} = 180\,$mm
\end{itemize}

\textbf{Động học thuận (FK).}
Cho trước $(\theta_1, \theta_2)$, tọa độ điểm cuối $P(x_P, y_P)$ xác định qua
ràng buộc vòng kín:
\begin{align}
    B_1 &= (L_{act}\cos\theta_1,\; L_{act}\sin\theta_1) \\
    B_2 &= (A_1A_2 + L_{act}\cos\theta_2,\; L_{act}\sin\theta_2)
\end{align}
$P$ là giao điểm của hai đường tròn bán kính $L_{pas}$ tâm tại $B_1$ và $B_2$
(chọn nghiệm $y_P < 0$ — chân hướng xuống).

\textbf{Động học nghịch (IK).}
Cho trước $P(x_P, y_P)$, tính ngược $(\theta_1, \theta_2)$:
\begin{align}
    \theta_1 &= \mathrm{atan2}(B_{1y},\; B_{1x}) \\
    \theta_2 &= \mathrm{atan2}(B_{2y},\; B_{2x} - A_1A_2)
\end{align}

\textbf{Ma trận Jacobian.}
Liên hệ vận tốc khớp với vận tốc điểm cuối:
\begin{equation}
    \begin{pmatrix} \dot{x}_P \\ \dot{y}_P \end{pmatrix}
    = J(\boldsymbol{\theta})
    \begin{pmatrix} \dot{\theta}_1 \\ \dot{\theta}_2 \end{pmatrix},
    \qquad J = \frac{\partial(x_P,\,y_P)}{\partial(\theta_1,\,\theta_2)}
\end{equation}
$J$ được tính số tại mỗi bước điều khiển; điểm kỳ dị ($\det J = 0$) ứng
với cấu hình cực biên không gian làm việc và cần tránh trong vận hành.

\endgroup
